% ID: SEL-D0055-OB-E02
% ESTADO: validada. S056.
Para $f$, escribimos $2^x/x+x$ y aplicamos el cociente usando $(2^x)'=2^x\ln2$:
\[
f'(x)=1+\frac{2^x(x\ln2-1)}{x^2},
\qquad x\ne0.
\]
En $g$ se aplican producto y cadena. El argumento del logaritmo es positivo para todo real:
\[
g'(x)=4x(x^2+1)\ln(e^{3x}+4)
+(x^2+1)^2\frac{3e^{3x}}{e^{3x}+4}.
\]
Para $h$, derivamos las dos potencias negativas y aplicamos la cadena en la segunda:
\[
h'(x)=-\frac{1}{3x^2}+\frac{10x}{(x^2-2)^2},
\qquad x\notin\{0,-\sqrt2,\sqrt2\}.
\]
